\section{Bad-history probability}
\label{sec:al_scan_bad_histories}

Keep the actual finite histories and physical graph from the preceding section.
Write $r=r_0$ and $M=\lceil C_1nD\rceil$. Time below counts completed
fill--charge pairs, so the charge at index $t$ sees the front after $t+1$ fills.
All probabilities in this section are classical finite-history probabilities.
The estimates formalize the finite ancestry argument of
\cite[Section~9, charge ancestry and availability]{OpenAI2026AreaLaw}; the
asymptotic scale specialization is proved below.

\begin{theorem}[Ancestry extracted from actual assignments]
    \label{thm:al_scan_actual_ancestry}
    \lean{TNLean.PEPS.AreaLaw.Scan.CollarScan.ChargeAncestry,
        TNLean.PEPS.AreaLaw.Scan.CollarScan.bandState_has_chargeAncestry,
        TNLean.PEPS.AreaLaw.Scan.CollarScan.oldChargeState_bad_implies_state_bad,
        TNLean.PEPS.AreaLaw.Scan.CollarScan.bandState_of_initial_assigned,
        TNLean.PEPS.AreaLaw.Scan.CollarScan.bad_bandState_initial_none,
        TNLean.PEPS.AreaLaw.Scan.CollarScan.ChargeAncestry.time_lt,
        TNLean.PEPS.AreaLaw.Scan.CollarScan.ChargeAncestry.ordered,
        TNLean.PEPS.AreaLaw.Scan.CollarScan.ChargeAncestry.selected,
        TNLean.PEPS.AreaLaw.Scan.CollarScan.ChargeAncestry.endpoint_mem_last_ball,
        TNLean.PEPS.AreaLaw.Scan.CollarScan.ChargeAncestry.anchor_chain,
        TNLean.PEPS.AreaLaw.Scan.CollarScan.bandState_has_chargeAncestry_domainGraph,
        TNLean.PEPS.AreaLaw.Scan.CollarScan.bad_bandState_mem_collar,
        TNLean.PEPS.AreaLaw.Scan.CollarScan.bad_bandState_has_long_chargeAncestry_domainGraph}
    \leanok
    \uses{def:al_scan_actual_histories}
    Suppose the target is nonempty, $8Km\le L$, and the actual cut endpoints
    have ambient distance greater than $2L+10r$ from the target.
    Every assigned site of $A$ has a chronological selected-charge path rooted
    at an initially assigned site or a scheduled fill slot. Its charge times
    are strictly increasing. For a nonempty path, its endpoint lies in its last charge ball, and
    successive anchors are at graph distance at most $2r$.
    Initial assignments persist. A bad endpoint was initially middle and lies
    in $A\cap T_L$. A scheduled fill cannot create a bad lead.
\end{theorem}
\begin{proof}\leanok
    Induct on the actual fill--charge evaluation. If the last charge makes a new
    assignment, choose a same-side predecessor in that ball. The ball meets the
    old middle, hence the cut clearance places the whole ball inside $A$.
    This includes far-side predecessors despite the exterior initially being far.
    Recurse on an earlier assignment, or terminate at the preceding fill.
    The resulting path family deliberately also permits selected charges that
    made no move; only containment of actual assignments in this family is used.
\end{proof}

\begin{theorem}[Length and time window of a bad ancestry]
    \label{thm:al_scan_ancestry_window}
    \lean{TNLean.PEPS.AreaLaw.Scan.CollarScan.ancestryFront,
        TNLean.PEPS.AreaLaw.Scan.CollarScan.front_mono,
        TNLean.PEPS.AreaLaw.Scan.CollarScan.front_speed,
        TNLean.PEPS.AreaLaw.Scan.CollarScan.ChargeAncestry.depth_le,
        TNLean.PEPS.AreaLaw.Scan.CollarScan.ChargeAncestry.bad_length_and_time}
    \leanok
    \uses{thm:al_scan_actual_ancestry}
    The fronts are monotone. For $n>0$ and pair times $a\le b$,
    \begin{align}
        b-a\le 2n\bigl(j(b)-j(a)+2\bigr).
    \end{align}
    If an assigned endpoint has lead greater than $D/2$, its ancestry is
    nonempty. For its length $s$ and every recorded charge time $t$,
    \begin{align}
        D<4rs,\qquad k-(t+1)\le 4n(rs+1).
    \end{align}
\end{theorem}
\begin{proof}\leanok
    Each charge changes oriented depth by at most $2r$. Keep the front at the
    first charge as the reference front, instead of replacing it by the final
    front. The exact padded fill counts give the displayed speed estimate with
    both floor errors included. Combining the two inequalities gives the bounds.
\end{proof}

\begin{theorem}[Probability of prescribed labelled hits]
    \label{thm:al_scan_selected_chain_probability}
    \lean{TNLean.PEPS.AreaLaw.Scan.CollarScan.sum_chargeWeight_selected_le,
        TNLean.PEPS.AreaLaw.Scan.CollarScan.sum_chargePathWeight_le_of_selected,
        TNLean.PEPS.AreaLaw.Scan.CollarScan.sum_chargePathWeight_le_of_selected_list,
        TNLean.PEPS.AreaLaw.Scan.CollarScan.sum_chargePathWeight_le_sum_of_selected_cover,
        TNLean.PEPS.AreaLaw.Scan.CollarScan.sum_chargePathWeight_le_sum_of_selected_list_cover}
    \leanok
    \uses{def:al_scan_actual_histories}
    Assume $M>0$ and fix all initial offsets. Requiring a specified side and label in one band
    at each of $s$ distinct times has probability at most $(2M)^{-s}$.
    The same upper bound holds after imposing arbitrary additional conditions.
    For a finite family of such prescriptions covering an event, the event's
    probability is at most the sum of these bounds.
\end{theorem}
\begin{proof}\leanok
    Sorted labelled candidate lists contain a label at most once, including when
    different labels share an anchor. Thus one required hit constrains at most
    one side--slot coordinate. Factor the finite product of independent choices
    and then take a union bound. Successful physical moves are not asserted
    independent; they are contained in the selected-coordinate events.
\end{proof}

\begin{theorem}[Counting spatial label paths]
    \label{thm:al_scan_anchor_path_count}
    \lean{TNLean.PEPS.AreaLaw.Scan.nearbyChargeLabels,
        TNLean.PEPS.AreaLaw.Scan.chargeAnchorPaths,
        TNLean.PEPS.AreaLaw.Scan.mem_chargeAnchorPaths_iff,
        TNLean.PEPS.AreaLaw.Scan.card_chargeAnchorPaths_le,
        TNLean.PEPS.AreaLaw.Scan.card_nearbyChargeLabels_le_of_vertex_count,
        TNLean.PEPS.AreaLaw.Scan.card_nearbyChargeLabels_domainGraph_le,
        TNLean.PEPS.AreaLaw.Scan.card_chargeAnchorPaths_domainGraph_le}
    \leanok
    \uses{thm:al_scan_actual_ancestry}
    Let $B$ bound the number of labels whose anchors lie within graph distance
    $2r$ of any site. The number of reversed label paths of length $s$ from a
    fixed endpoint is at most $B^s$. In an induced square-lattice domain with
    anchor multiplicity at most $\mu$, one may take
    \begin{align}
        B=\mu\bigl(1+2(2r)(2r+1)\bigr).
    \end{align}
\end{theorem}
\begin{proof}\leanok
    Recursively choose a nearby labelled anchor and continue from it. The finite
    union has at most $B$ choices at each step. Fiber counting retains label
    multiplicity, and an induced graph ball injects into the ambient lattice
    diamond. Counting labels alone suffices because predecessor witnesses do
    not enter the selected-coordinate probability.
\end{proof}

\begin{theorem}[Distinct charge times in a terminal interval]
    \label{thm:al_scan_timed_path_count}
    \lean{TNLean.PEPS.AreaLaw.Scan.timedChargePaths,
        TNLean.PEPS.AreaLaw.Scan.card_timedChargePaths_le,
        TNLean.PEPS.AreaLaw.Scan.mem_timedChargePaths,
        TNLean.PEPS.AreaLaw.Scan.length_of_mem_timedChargePaths,
        TNLean.PEPS.AreaLaw.Scan.time_mem_of_mem_timedChargePaths,
        TNLean.PEPS.AreaLaw.Scan.ordered_of_mem_timedChargePaths,
        TNLean.PEPS.AreaLaw.Scan.recentChargeTimes,
        TNLean.PEPS.AreaLaw.Scan.card_recentChargeTimes_le,
        TNLean.PEPS.AreaLaw.Scan.mem_recentChargeTimes,
        TNLean.PEPS.AreaLaw.Scan.card_ancestry_time_window_le}
    \leanok
    \uses{thm:al_scan_ancestry_window}
    A terminal window specified by $k-(t+1)\le W$ and $0\le t<k$ contains
    at most $W+1$ indices. Choosing $s$ of its $v$ indices and a reversed label
    path from a finite family of size $b$ gives at most $\binom vs b$
    chronologically labelled paths. Their times are distinct.
    For $n,r,s>0$ and $W=4n(rs+1)$, one has $v\le 9nrs$.
\end{theorem}
\begin{proof}\leanok
    Use the subsets of cardinality $s$, sort them, and pair their times with the
    label path in chronological order. The image cannot increase cardinality.
    Keeping the inclusive endpoint produces $W+1$; positivity absorbs it and
    the additive rounding terms into $9nrs$.
\end{proof}

\begin{theorem}[Finite cover of actual bad ancestries]
    \label{thm:al_scan_bad_path_cover}
    \lean{TNLean.PEPS.AreaLaw.Scan.CollarScan.ChargeAncestry.labels_mem_chargeAnchorPaths,
        TNLean.PEPS.AreaLaw.Scan.CollarScan.ChargeAncestry.length_le,
        TNLean.PEPS.AreaLaw.Scan.CollarScan.badChargePathsOfLength,
        TNLean.PEPS.AreaLaw.Scan.CollarScan.badChargePaths,
        TNLean.PEPS.AreaLaw.Scan.CollarScan.ChargeAncestry.mem_badChargePaths,
        TNLean.PEPS.AreaLaw.Scan.CollarScan.length_of_mem_badChargePathsOfLength,
        TNLean.PEPS.AreaLaw.Scan.CollarScan.time_lt_of_mem_badChargePaths,
        TNLean.PEPS.AreaLaw.Scan.CollarScan.distinct_times_of_mem_badChargePaths,
        TNLean.PEPS.AreaLaw.Scan.CollarScan.card_badChargePathsOfLength_le,
        TNLean.PEPS.AreaLaw.Scan.CollarScan.sum_badChargePaths_le}
    \leanok
    \uses{thm:al_scan_actual_ancestry,thm:al_scan_ancestry_window,thm:al_scan_anchor_path_count,thm:al_scan_timed_path_count}
    For a fixed endpoint and pair time $k$, all bad ancestries lie in the finite
    union of the timed spatial path families with $0\le s\le k$ and $D<4rs$.
    Every path in this family has valid distinct times. For every $q\ge0$, its
    sum of weights $q^s$ is at most
    \begin{align}
        \sum_{\substack{0\le s\le k \\D<4rs}}
        \binom{v_s(k)}s(Bq)^s,
        \qquad v_s(k)=\min\{k,4n(rs+1)+1\}.
    \end{align}
\end{theorem}
\begin{proof}\leanok
    The reversed labels satisfy the spatial chain relation, and strict time order
    bounds the path length by $k$. Insert the actual path into its sorted-time
    family. Families of different lengths are disjoint, so summing their
    cardinality bounds gives the displayed estimate.
\end{proof}

\begin{theorem}[Binomial cancellation]
    \label{thm:al_scan_binomial_cancellation}
    \lean{TNLean.PEPS.AreaLaw.Scan.choose_le_exp_mul_div_pow,
        TNLean.PEPS.AreaLaw.Scan.choose_mul_pow_le_of_window_bound}
    \leanok
    \uses{thm:al_scan_timed_path_count}
    For $s>0$, $a\ge0$, and $v\le Cs$,
    \begin{align}
        \binom vs\le (ev/s)^s,\qquad
        \binom vs a^s\le(eCa)^s.
    \end{align}
\end{theorem}
\begin{proof}\leanok
    Use $\binom vs\le v^s/s!$ and $s^s/s!\le e^s$. The factorial cancels
    the ancestry length in the time-window estimate.
\end{proof}

\begin{theorem}[Compact collar cardinality]
    \label{thm:al_scan_compact_collar_card}
    \lean{TNLean.PEPS.AreaLaw.Scan.card_ambientDilation_le_of_layers,
        TNLean.PEPS.AreaLaw.Scan.card_compact_collar_le}
    \leanok
    \uses{def:al_scan_actual_histories}
    If every positive ambient layer through $L$ has at most $n$ sites, then
    \begin{align}
        |A\cap T_L|\le |T_L|\le |T|+nL.
    \end{align}
\end{theorem}
\begin{proof}\leanok
    The zero dilation is $T$. Add the cardinalities of successive actual dilation
    differences and inject physical collar sites into the ambient dilation.
    No volume of $\Lambda\setminus A$ occurs.
\end{proof}

\begin{theorem}[Exact prefix-history marginals]
    \label{thm:al_scan_prefix_marginal}
    \lean{TNLean.PEPS.AreaLaw.Scan.prefixHistory,
        TNLean.PEPS.AreaLaw.Scan.prefixHistory_self,
        TNLean.PEPS.AreaLaw.Scan.prefixHistory_extendHistory,
        TNLean.PEPS.AreaLaw.Scan.sum_historyWeight_mul_prefix}
    \leanok
    \uses{def:al_scan_actual_histories}
    Restricting a full history to its first $j$ pairs preserves its offsets and
    earlier charge coordinates. For $j\le N$ and $M>0$, every real-valued
    function of this prefix has the same expectation under length-$N$ history
    weights as under length-$j$ history weights.
\end{theorem}
\begin{proof}\leanok
    Appending a charge preserves all earlier prefixes. Sum the appended independent
    choices, whose total weight is one, and induct from $j$ to $N$. Empty
    offset spaces give zero on both sides and require no extra convention.
\end{proof}

\begin{theorem}[Finite bad-endpoint bound]
    \label{thm:al_scan_bad_endpoint_probability}
    \lean{TNLean.PEPS.AreaLaw.Scan.CollarScan.IsBadBandEndpoint,
        TNLean.PEPS.AreaLaw.Scan.CollarScan.badEndpointTail,
        TNLean.PEPS.AreaLaw.Scan.CollarScan.sum_chargePathWeight_bad_endpoint_domainGraph_le}
    \leanok
    \uses{thm:al_scan_bad_path_cover,thm:al_scan_selected_chain_probability,thm:al_scan_anchor_path_count}
    Fix the offsets, a band, a side, and an endpoint in $A$. Under the actual
    cut-clearance hypotheses and $n,M>0$, the probability that the endpoint is
    assigned with lead greater than $D/2$ is at most
    \begin{align}
        R_k=\sum_{\substack{0\le s\le k \\D<4rs}}
        \binom{v_s(k)}s\left(\frac{B}{2M}\right)^s.
    \end{align}
\end{theorem}
\begin{proof}\leanok
    Derive the ancestry from the evaluated physical history, put it in the finite
    cover, and apply the selected-coordinate probability bound to each path.
    The graph and ambient depth provide the spatial variation and diamond count.
\end{proof}

\begin{theorem}[Bands, sides, compact endpoints, and offsets]
    \label{thm:al_scan_bad_history_probability}
    \lean{TNLean.PEPS.AreaLaw.Scan.CollarScan.IsGoodCompletedHistory,
        TNLean.PEPS.AreaLaw.Scan.CollarScan.isGoodOldHistory_of_isGoodCompletedHistory,
        TNLean.PEPS.AreaLaw.Scan.CollarScan.sum_chargePathWeight_bad_completed_domainGraph_le,
        TNLean.PEPS.AreaLaw.Scan.CollarScan.sum_historyWeight_bad_completed_domainGraph_le,
        TNLean.PEPS.AreaLaw.Scan.CollarScan.sum_historyWeight_bad_old_domainGraph_le}
    \leanok
    \uses{thm:al_scan_bad_endpoint_probability,thm:al_scan_actual_ancestry}
    For $m>0$, the total weight of bad completed histories at time $k$ is at most
    $2K|A\cap T_L|R_k$. The same bound holds for old charge histories after
    their preceding fill.
\end{theorem}
\begin{proof}\leanok
    Every bad endpoint belongs to the compact collar. Union over it, both sides,
    and all bands, then average the normalized initial offsets. A fill cannot
    create a bad lead, so old charge states are covered as well.
\end{proof}

\begin{theorem}[All actual times on one full history]
    \label{thm:al_scan_bad_all_times}
    \lean{TNLean.PEPS.AreaLaw.Scan.CollarScan.IsGoodThrough,
        TNLean.PEPS.AreaLaw.Scan.CollarScan.IsGoodThrough.old,
        TNLean.PEPS.AreaLaw.Scan.CollarScan.sum_historyWeight_bad_through_domainGraph_le}
    \leanok
    \uses{thm:al_scan_bad_history_probability,thm:al_scan_prefix_marginal,thm:al_scan_compact_collar_card}
    Under the source ambient layer bound, the probability that some completed
    prefix of one length-$N$ actual history is bad is at most
    \begin{align}
        2K(|T|+nL)\sum_{k=0}^N R_k.
    \end{align}
    Goodness of these completed prefixes also guarantees goodness after every
    intervening deterministic fill.
\end{theorem}
\begin{proof}\leanok
    Take the finite union over prefix times and use their exact marginals. Insert
    the compact collar cardinality bound. All tests concern the same full
    history, rather than independently sampled histories at different times.
\end{proof}

\begin{theorem}[Finite geometric bad-history tail]
    \label{thm:al_scan_bad_geometric_tail}
    \lean{TNLean.PEPS.AreaLaw.Scan.CollarScan.badChargeRatio,
        TNLean.PEPS.AreaLaw.Scan.CollarScan.geometricBadTail,
        TNLean.PEPS.AreaLaw.Scan.CollarScan.badEndpointTail_le_geometricBadTail,
        TNLean.PEPS.AreaLaw.Scan.CollarScan.geometricBadTail_mono,
        TNLean.PEPS.AreaLaw.Scan.CollarScan.sum_historyWeight_bad_through_geometric_domainGraph_le}
    \leanok
    \uses{thm:al_scan_bad_all_times,thm:al_scan_binomial_cancellation}
    If additionally $r>0$, set
    \begin{align}
        \rho=\frac{9e\,nrB}{2M}.
    \end{align}
    The probability of a bad prefix is at most
    \begin{align}
        2K(|T|+nL)(N+1)
        \sum_{\substack{0\le s\le N \\D<4rs}}\rho^s.
    \end{align}
\end{theorem}
\begin{proof}\leanok
    Apply the binomial cancellation with $v_s(k)\le9nrs$. The nonnegative
    length sum increases with the observation horizon. The label-only count
    gives a ratio of order $r^3/D$ when $M$ has the source scale $nD$,
    with uniformly bounded anchor multiplicity and fixed positive $C_1$;
    counting predecessor sites as well gives the manuscript's looser $r^5/D$.
    No asymptotic limit or transported quantum measure is used here.
\end{proof}

\begin{figure}[ht]
    \centering
    % Physical-site diagram, not a tensor contraction: no wires or open ports.
    % Columns are ambient depths 32..39 in the active path component.
    % Populated rows follow the three radius-one charges; the front stays 33.
    % Empty wire rows reserve space for the north-facing charge labels.
    % Enclosures have members 33..35, 35..37, 37..39, respectively.
    \tenkzkernel
    \begin{tenkz}[rows={wire,wire,wire,wire,wire}, cols=8, bonds=none, west=none, east=none]
        \tn[skin=box]{32} & \tn[skin=box, name=scanFirstRow33]{33} &
        \tn[skin=box]{34} & \tn[skin=box]{35} &
        \tn[skin=box]{36} & \tn[skin=box]{37} &
        \tn[skin=box]{38} & \tn[skin=box]{39} \\
        \\
        \tn[skin=box]{32} & \tn[skin=box]{33} &
        \tn[skin=box]{34} & \tn[skin=box]{35} &
        \tn[skin=box]{36} & \tn[skin=box]{37} &
        \tn[skin=box]{38} & \tn[skin=box]{39} \\
        \\
        \tn[skin=box]{32} & \tn[skin=box, name=scanFront]{33} &
        \tn[skin=box]{34} & \tn[skin=box]{35} &
        \tn[skin=box]{36} & \tn[skin=box, name=scanGoodThreshold]{37} &
        \tn[skin=box]{38} & \tn[skin=box, name=scanBadEndpoint]{39}
        \tnmark[form=enclosure, species=selected, tint, name=scanCharge34,
            label pos=n]{(1,2)..(1,4)}{$B_1(34)$}
        \tnmark[form=enclosure, species=selected, tint, name=scanCharge36,
            label pos=n]{(3,4)..(3,6)}{$B_1(36)$}
        \tnmark[form=enclosure, species=selected, tint, name=scanCharge38,
            label pos=n]{(5,6)..(5,8)}{$B_1(38)$}
        \tnmark[form=label, label pos=s]{scanFront}{$j=33$}
        \tnmark[form=label, label pos=s]{scanGoodThreshold}{$j+D/2$}
        \tnmark[form=label, label pos=s]{scanBadEndpoint}{$39>37$}
    \end{tenkz}
    \caption{Three actual charges, from top to bottom, with $m=32$, $D=8$,
        and $r_0=1$. The preceding fill assigns site $33$; the successive charge balls
        extend the near side through sites $35$, $37$, and $39$ while the nominal front remains $33$.
        Their overlaps supply the ancestry predecessors. The displayed sites
        form the induced path component $32$--$39$. Ambient depth is measured from
        the separate target component $T=\{(0,0)\}$, whereas each charge ball
        uses graph distance inside the active component. A separate physical
        cut $2000$--$2001$ provides genuine cut clearance; neither distant component
        is joined to the displayed path.}
    \label{fig:al_scan_three_charge_ancestry}
\end{figure}

\begin{lemma}[Rounded logarithmic radius]
    \label{lem:al_scan_rounded_radius}
    \lean{TNLean.PEPS.AreaLaw.Scan.roundedLogRadius,
        TNLean.PEPS.AreaLaw.Scan.isLittleO_roundedLogRadius_rpow,
        TNLean.PEPS.AreaLaw.Scan.eventually_const_mul_roundedLogRadius_le_rpow,
        TNLean.PEPS.AreaLaw.Scan.eventually_roundedLogRadius_le_rpow,
        TNLean.PEPS.AreaLaw.Scan.eventually_const_mul_roundedLogRadius_cube_le_rpow,
        TNLean.PEPS.AreaLaw.Scan.eventually_one_le_roundedLogRadius}
    \leanok
    Fix $C_r\ge0$ and $r(n)=\lceil C_r(\log n)^2\rceil$.
    For every $\varepsilon>0$, one has $r(n)=o(n^\varepsilon)$.
    Consequently, for every fixed real $A$, both
    $Ar(n)\le n^\varepsilon$ and $Ar(n)^3\le n^\varepsilon$
    hold for sufficiently large $n$. If $C_r>0$, then $r(n)\ge1$
    for sufficiently large $n$.
\end{lemma}
\begin{proof}\leanok
    The ceiling differs from the nonnegative squared logarithm by less than one.
    Every fixed power of the logarithm is negligible relative to every positive
    power of $n$. Apply this first with exponent $\varepsilon$ and then with
    exponent $\varepsilon/3$ before cubing.
\end{proof}

\begin{lemma}[A fixed ancestry cutoff]
    \label{lem:al_scan_fixed_ancestry_cutoff}
    \lean{TNLean.PEPS.AreaLaw.Scan.CollarScan.badChargeRatio_le_radius_cube,
        TNLean.PEPS.AreaLaw.Scan.CollarScan.source_lengths_le,
        TNLean.PEPS.AreaLaw.Scan.CollarScan.geometricBadTail_le_power,
        TNLean.PEPS.AreaLaw.Scan.CollarScan.badHistory_prefactor_le,
        TNLean.PEPS.AreaLaw.Scan.CollarScan.badHistory_power_le}
    \leanok
    \uses{thm:al_scan_bad_geometric_tail}
    Write $b$ for the fixed anchor multiplicity bound. For $n,D,r>0$ and
    $M\ge C_1nD$, where $C_1>0$, the actual ratio satisfies
    \begin{align}
        0\le\rho\le \frac{117e b}{2C_1}\frac{r^3}{D}.
    \end{align}
    If $n\ge1$, $a\ge0$, $4rJ\le D$, and $\rho\le n^{-a}$, then
    \begin{align}
        \sum_{\substack{0\le s\le N \\D<4rs}}\rho^s
        \le (N+1)n^{-aJ}.
    \end{align}
    If moreover $K,L,m\le n$, $|T|\le C_Tn^2$, $C_T\ge0$, and $N=nm$, then
    \begin{align}
        2K(|T|+nL)(N+1)^2\le8(C_T+1)n^7.
    \end{align}
    For $n\ge8(C_T+1)$ and $p+8\le aJ$, this prefactor times
    $n^{-aJ}$ is at most $n^{-p}$.
\end{lemma}
\begin{proof}\leanok
    For $r\ge1$ the lattice diamond factor is at most $13r^2$.
    Cancel $n$ against the padded capacity. Every retained length is at least
    $J$, and there are at most $N+1$ lengths. Bound $N+1$ by $2n^2$ and absorb
    the remaining fixed constant into one additional power of $n$.
\end{proof}

\begin{theorem}[Bad-history probability at the source scales]
    \label{thm:al_scan_bad_history_inverse_power}
    \lean{TNLean.PEPS.AreaLaw.Scan.eventually_badHistory_bound_le_rpow,
        TNLean.PEPS.AreaLaw.Scan.eventually_bad_history_probability_le}
    \leanok
    \uses{thm:al_scan_bad_geometric_tail,lem:al_scan_rounded_radius,
        lem:al_scan_fixed_ancestry_cutoff}
    Fix $0<\ell$, $0<\kappa<\mu<1-\ell$, $C_r,C_1>0$, $C_T\ge0$,
    a finite anchor multiplicity bound $b$, and a real exponent $p$.
    Set
    \begin{align}
        r & =\lceil C_r(\log n)^2\rceil, & D                      & =\lceil n^\kappa\rceil,
          & m                            & =\lfloor n^\mu\rfloor,                           \\
        L & =\lfloor n^{1-\ell}\rfloor,  & K                      & =\lfloor L/(8m)\rfloor,
          & M                            & =\lceil C_1nD\rceil.
    \end{align}
    For all sufficiently large $n$, uniformly over finite lattice domains,
    nonempty targets with $|T|\le C_Tn^2$, and actual scanners satisfying
    the ambient row bound, the cut clearance above, and anchor multiplicity
    at most $b$, the probability that any prefix through $nm$ fill--charge
    pairs has bad lead on either side in any band is at most $n^{-p}$.
    In particular $p=200$ gives the bad-history exponent in Lemma~9.1.
\end{theorem}
\begin{proof}\leanok
    Choose an integer $J$ with $p+8\le\kappa J/2$, before choosing $n$ or
    the physical domain. The rounded-radius estimates imply eventually
    $\rho\le n^{-\kappa/2}$ and $4rJ\le D$. The source floors give
    $K,L,m\le n$. Apply the fixed-cutoff estimate to the actual finite-history
    union bound. The same good-prefix event controls the old-charge states,
    since a scheduled fill cannot create a bad lead.
\end{proof}

Identifying transported entropy gains and constructing the remaining scanner
inputs are separate mathematical steps. The probability estimate concerns
actual classical histories and does not identify their quantum entropy costs.
